\documentclass[draft]{agujournal2025}
\usepackage{booktabs}
\usepackage{amsmath}
\usepackage{amssymb}
\usepackage{tabularx}
\usepackage{microtype}
\usepackage{url}
\setlength{\emergencystretch}{3em}
\newsavebox{\fitbox}
\newcommand{\fittable}[1]{\sbox{\fitbox}{#1}\ifdim\wd\fitbox>\textwidth\resizebox{\textwidth}{!}{\usebox{\fitbox}}\else\usebox{\fitbox}\fi}
\input{tr_claims.tex}
\newcommand{\Hthirty}{H30}
\begin{document}
\journalname{Water Resources Research}

\title{Incomplete river-temperature years are more useful than completeness rules assume}

% Author, affiliation and funding information must be supplied by the authors; nothing is invented here.
\authors{AUTHOR NAMES PENDING AUTHOR CONFIRMATION\affil{1}}
\affiliation{1}{AFFILIATION PENDING AUTHOR CONFIRMATION}
\authoraddr{Corresponding author details pending author confirmation.}
\authorrunninghead{Ignored.}
\titlerunninghead{Ignored.}

\keypoints
  {Gap filling understates the annual thermal peak; a recovery law predicts where, but underpredicts failure in the hardest cases}
  {Completeness rules discard mostly recoverable years; keeping them gives most of the gain in station trend accuracy}
  {Fixed admission thresholds failed in a harder region; a label-free threshold based on expected risk still needs an independent test}

\maketitle\par

\begin{abstract}
River-temperature archives contain many incomplete years, which analysts either discard under completeness rules or fill before computing annual extremes. We tested both practices for the hottest 30-day mean (H30) on \TSourceYears{} complete station-years from \TSourceStations{} stations, with real gap calendars transplanted onto complete years, under two preregistrations. Missing data were only weakly related to heat. Filled series underestimated H30 when its window was hidden, by \TPOneBiasPZeroAbs\,$^{\circ}$C where only air temperature informed the gap; a distribution-regression estimator removed \TPOneReductionCal\% of this bias, with approximately calibrated intervals; the registered endpoint was missed narrowly, on interval coverage (\TPOneCoverageRone\% against a floor of 85\%). A recovery law built from training data alone explained \TLawCellFailEV\% of the variation in failure rates across gap geometries but underpredicted the hardest cases. Published completeness rules used \TRetainMin--\TRetainMax\% of incomplete years, and \TRejectedUsableMin--\TRejectedUsableMax\% of the rejected years were recoverable within 0.5\,$^{\circ}$C. Keeping all years through multiple imputation reduced the error of station H30 trends by \TTrendReduction\% relative to a 90\% warm-season rule (preregistered endpoint met); in a post hoc decomposition, \TShareKeep\% of this gain came from keeping the years. Fixed admission thresholds failed in Australia. A label-free threshold developed afterwards kept the admitted risk at \TNewRiskDeploy\% and \TNewRiskMech\% at \TNewStations{} new USGS stations, within the preregistered 7\%; both tests were recomputed after their first outcomes were known and are no longer blind. In an exploratory test, the same mechanism biased annual 7-day minimum flows upward. Archive users should estimate the statistic for incomplete years and carry its uncertainty forward rather than discard them.
\begin{plainlanguagesummary}
Long river-temperature records are used to measure warming and heat waves, but many years have gaps. Common practice either drops those years or fills the gaps and then reads off the warmest month. We tested both choices on thousands of complete years in which we hid days on purpose. Filling the gaps made hot summers look cooler than they were, especially when only air temperature was available to fill them. A method that estimates the warmest-month value directly, with an approximate uncertainty range, removed most of this error. Dropping years was wasteful: most years thrown out by published rules could have been estimated accurately, and keeping them made station warming trends several times more accurate. A rule for deciding which estimates are good enough failed in Australia when it used fixed thresholds; a revised rule worked at new river stations, but we recomputed that check after seeing its first result, so it still needs an independent test. Low river flows showed the same problem in reverse: filled records made dry spells look wetter. Gaps do not mainly hide the hottest years, but how analysts handle them changes conclusions.
\end{plainlanguagesummary}
\end{abstract}

\section{Introduction}
River-temperature records underpin assessments of warming, heat waves and habitat suitability \cite{caissie2006,michel2020,vanhamel2024}. Many of these assessments use annual statistics of sustained heat, such as the mean of the hottest 30 consecutive days (H30), which \citeA{beaufort2022} introduced as the stream thermal peak. A large share of station-years in operational archives are incomplete. In the archive used here, \TWarmIncomplete\% of fully operational station-years have at least one missing day between May and September.

Practitioners handle incomplete years in two ways. The first is to discard years that fail a completeness rule. Rules that count missing days include the USGS statistics services, which do not compute annual statistics when any day is missing \cite{kiang2018}, the USGS stream-temperature indicators \cite{johnson2024usgs}, the WMO 11/5 rule for monthly values \cite{wmo2017}, the ETCCDI rule \cite{zhang2011} and the 350-day rule of GSIM \cite{gudmundsson2018}. Other rules look at the season in which the peak occurs: Oregon's assessment methodology requires 80\% of days from July to September \cite{oregon2024}, NorWeST uses August records with at least 90\% coverage \cite{isaak2017}, and \citeA{beaufort2022} accept years with complete July and August records. The second practice is to fill the gaps with a regression on air temperature or neighbouring stations, or with a process model, and to compute the statistic from the completed series \cite{toffolon2015,piccolroaz2016}.

Both practices can bias thermal extremes. A rule that discards years removes information, and if gaps are more frequent in hot, low-flow summers it also removes the hottest years \cite{sowder2012}. Filling and then taking the maximum is biased for a different reason. A regression returns the conditional mean, so it shrinks hidden anomalies towards the seasonal cycle, and the maximum of a smoothed series is lower than the maximum of the true one. Statistical texts warn that imputation is not prediction and that the imputation with the lowest squared error is a poor basis for estimating features of a distribution \cite{vanbuuren2018,naf2023}. Air-temperature regressions of stream temperature are known to miss the warmest and coldest periods \cite{arismendi2014}. \citeA{johnson2021} showed with masking experiments that the tolerance of annual stream-temperature parameters to gaps depends on gap length and on the parameter, but they did not study maxima and noted that gaps near the annual maximum matter most.

Three things are missing. There is no estimate of how often real gaps coincide with hot states. There is no estimator designed for the statistic rather than for the daily series, together with an assessment of whether it beats current rules at equal risk. And there is no test of whether the choice changes the scientific conclusions that archives are used for, such as station warming trends. Machine-learning work on missing data has shown that the value of an imputation depends on the downstream task \cite{lemorvan2021,shadbahr2023,wang2024imputation}, but this principle has not been tested on hydrological extremes.

We address four questions. Is missingness in river-temperature archives related to hot or low-flow states? Can the annual extreme be recovered without bias and with approximately calibrated intervals, and can the error of the current practice be predicted from training data alone? How many incomplete years do current rules discard that could have been used, and does a target-aware applicability judgment do better at equal risk? Does the treatment of incomplete years change station trends and hot-year counts? Two preregistered tests on data that were not used in development ask whether an admission threshold transfers, and a demonstration on annual 7-day minimum flows asks whether the mechanism is specific to temperature.

\section{Data and Methods}

\subsection{Archive and hot years}
The study uses daily water temperature from ten agencies in Europe and the United States that were assembled for an earlier study and have already been analysed. The full archive comprises \TArchiveStations{} stations and \TArchiveYears{} station-years, complete or not. Daily air temperature for each station is the NASA POWER 2-m product \cite{nasapower2026}, which we extended to 132 USGS stations that had no forcing. Daily discharge was retrieved for 287 USGS stations \cite{usgs2026}. Stations are grouped into independence groups of record-connected networks, which are the units of resampling, and into regions for transfer tests: each European agency forms a region and USGS stations are split by two-digit hydrologic unit.

Annual statistics are computed in hot years. For Northern Hemisphere stations a hot year is the calendar year; Southern Hemisphere records are shifted by 184 days so that a hot year runs from July to June and the summer peak lies away from the year boundary. The warm season of a station is its five climatologically hottest consecutive months, taken from the training years of each case, and degree-days are summed over its four hottest months. The first version of this analysis used calendar years and a May--September season everywhere; the correction leaves the source masks, H30 values and plug-in values unchanged, but \TDefWarmNotMaySep{} of \TDefStations{} source stations have a different warm season (Table~S18).

\subsection{Missingness audit}
We analysed all fully operational station-years from 1982 onwards, excluding the first and last year of each record (\TEOneYears{} station-years at \TEOneStations{} stations). At the station-year level, a conditional logistic regression with station fixed effects related warm-season incompleteness to the within-station standardized summer air-temperature anomaly and a decadal trend; a mixed model with station and region-year random intercepts served as a sensitivity analysis. At the daily level, a time-stratified case-crossover design \cite{maclure1991,janes2005} compared the 7-day air-temperature anomaly, the lagged water-temperature anomaly and, for USGS stations, a low-flow indicator on gap-onset days with the other days at risk in the same station-month. Intervals resample independence groups.

\subsection{Case populations}
Every complete station-year was masked repeatedly, with the target year excluded from all fitting. The mechanism population contains one contiguous 30-, 60- or 90-day gap with no, partial or full overlap with the training-period climatological peak window, two replicates per cell (\TCrossedCases{} cases). The deployment population transplants the missing-day calendar of a randomly drawn real incomplete year of the same agency onto the complete year, shifting each gap by up to 15 days (\TCalendarCases{} cases); it is weighted so that the agency mix matches the real incomplete years. State-dependent missingness is represented by exact importance weights $\prod_j \exp(k\beta z_{t_j}) / \overline{\exp(k\beta z)}_{w_j}$ on the calendar cases, where $z_{t_j}$ is the standardized 7-day air anomaly on the shifted onset day of gap $j$, $\beta$ is the onset effect estimated in the audit, the average runs over the 31-day shift window $w_j$, and $k\in\{1,2\}$. In total, \TSourceCases{} cases come from \TSourceYears{} station-years at \TSourceStations{} stations in \TSourceGroups{} independence groups and \TSourceRegions{} regions.

Gaps are filled by an information-priority policy that uses neighbour-station ridge regression when at least one neighbour covers 90\% of the hidden days, else air-temperature ridge regression with lags of up to 14 days, else linear interpolation for gaps up to 30 days, else a smoothed day-of-year climatology. The information layer is the method chosen; \TAirLayerShare\% of cases use the air-temperature layer.

\subsection{Targets and tolerances}
The primary target is H30 with a tolerance of 0.5\,$^{\circ}$C, the accuracy limit of a ``good'' continuous temperature record \cite{wagner2006} and about half the median interannual standard deviation of H30. The other targets are the hottest 7-day mean (H7), the annual mean, the coldest 30-day mean (C30), the degree-days of the four hottest months and the number of days above 20\,$^{\circ}$C. Their tolerances were set to the same fraction of their interannual variability: \TTolHSeven, \TTolMean{} and \TTolCThirty\,$^{\circ}$C, \TTolDD{} degree-days and \TTolDays{} days. A value is a failure when its error exceeds the tolerance.

\subsection{Estimators}
The plug-in estimator P0 computes the target from the completed series; it was the value delivered by earlier work. We compare it with the target computed from recorded days only, requiring 80\% of days in a window (the most common practice); Air2stream-3 \cite{toffolon2015} calibrated on the longest record segment outside the target year and started from the last observation before each gap (P1); functional debiasing, which subtracts a cross-fitted gradient-boosting \cite{ke2017} prediction of the signed error of P0 (P2); the median of 128 conditional simulations that add seasonal residual blocks and slow residual analogues selected by the target year's air-temperature state to the reconstruction, with raw residual blocks as an ablation (P3); and direct distribution regression of H30, which predicts seven quantiles of the H30 error with gradient boosting, without (P5a) or with (P5b) the simulation summaries as inputs. All learned components were cross-fitted by leaving one region out. Following the preregistration, the delivered distribution D2* is the candidate with the lowest mean quantile continuous ranked probability score \cite{gneiting2007}; squared error was not a selection criterion \cite{naf2023}.

The quantiles of D2* were then calibrated by split-conformal offsets for each quantile level \cite{romano2019}, estimated leave-one-region-out within groups that are known at decision time: the share of the climatological peak window that is hidden (none, or one of four quarters) crossed with the information layer, with small layers merged. This step was added after the first submission, when the 90\% intervals were found to under-cover where most of the peak was hidden.

\subsection{Recovery law}
Write the daily water temperature as a seasonal cycle plus an anomaly $a_t$ with warm-season standard deviation $\sigma_a$ and lag-one correlation $\phi_a$, and the reconstruction residual (centre minus truth) as $r_t$ with mean $\mu_r$, standard deviation $\sigma_r$ and lag-one correlation $\phi_r$. If a conditional-mean reconstruction leaves the residual uncorrelated with its prediction, then $\mathrm{Cov}(r_t, a_s) = -\mathrm{Cov}(r_t, r_s)$. The peak is selected on the 30-day mean, so the relevant regression is that of the hidden residual mean $\bar r_m$ on the window anomaly $\bar a$:
\begin{equation}
E[\bar r_m \mid \bar a] = \mu_r - \kappa_m \bar a, \qquad \kappa_m = \frac{\sigma_r^2\, c(m,\phi_r)}{\sigma_a^2\, v(30,\phi_a)},
\end{equation}
where $v(m,\phi) = [m + 2\sum_{j=1}^{m-1}(m-j)\phi^j]/m^2$ is the variance factor of an $m$-day mean of an AR(1) process and $c(m,\phi) = (30m)^{-1}\sum_{t\in\mathcal H}\sum_{s\in\mathcal S}\phi^{|t-s|}$ couples the $m$ hidden days $\mathcal H$ with the 30-day window $\mathcal S$. If a share $h=m/30$ of the window is hidden, the plug-in error is approximately
\begin{equation}
e \sim \mathcal N\!\left(h(\mu_r - \kappa_m \bar a),\; h^2\left[\sigma_r^2 v(m,\phi_r) - \frac{\sigma_r^4 c(m,\phi_r)^2}{\sigma_a^2 v(30,\phi_a)}\right]\right).
\end{equation}
The position and anomaly of the target year's peak are unknown, so the law mixes over the training years' own pairs $(h_y, \bar a_y)$, obtained by placing the current gap on each training year's true peak window. Because H30 cannot be lower than the best fully observed window $L$, the mixture is truncated at $e \le \mathrm{P0} - L$. All quantities come from the training years; nothing is fitted to reconstruction errors. We also report a daily-scale version that uses $1-R^2$ instead of $\kappa_m$ (v1) and a version that takes the residual moments inside the training peak windows (v2p). Replaying the gap on each training year with that year's cross-fitted residuals (v3) gives an empirical comparator; it is not a law.

\subsection{Applicability judgments and published rules}
We implemented fifteen published rules (Table~\ref{tab:rules}). Day-counting rules are applied as written. Peak-season rules were written for July--September and its parts in the Northern Hemisphere; we anchor their windows to each station's climatological peak, which adapts them to other climates, and also report them as written (Table~\ref{tab:rules}). The July--September rule follows Oregon's 80\% coverage requirement using daily means rather than 7-day average daily maxima. Three rules with a free threshold were also traced as curves: annual missing fraction, warm-season missing fraction, and an NRFA-type observed lower bound \cite{nrfa2026}. Learned or derived scores were the daily and H30 disagreement among reconstructions, the recovery-law failure probability, the D2* predictive failure probability $P(|\mathrm{H30}-\hat{\mathrm{H30}}|>0.5)$ (R7a), and a gradient-boosting classifier of the failure of the D2* value (R7b), which we call the target-aware judgment. A calibrated failure probability is the optimal score for selecting estimates with a fixed point \cite{chow1970,franc2023}. Scores and rules were compared on risk--retention curves at equal realised risk, and by AURC and AUGRC \cite{geifman2017,traub2024}; retention is the share of incomplete years used and risk the share of used values outside tolerance. Each rule was combined with each delivered value (P0, recorded days or D2*).

\subsection{Transfer and the two preregistered tests}
\label{sec:tests}
The first preregistration fixed admission thresholds in probability units: for a nominal 5\% admitted risk, thresholds were set on the other regions' out-of-fold scores either as one threshold or as one threshold per information layer \cite{vovk2005}, also with covariate-shift weights \cite{tibshirani2019} and with local calibration blocks. After the source analyses, the pipeline (estimator selection, models fitted on all source cases, thresholds per layer from the source deployment population, and the code) was frozen, and only then were cases generated for the Australian Bureau of Meteorology's qualified daily series \cite{bom2026}: \TAuStations{} stations with at least four complete years in \TAuNetworks{} hydrological networks, \TAuYears{} station-years and \TAuCases{} cases. This first test was evaluated once at the time.

After it failed, we developed a label-free rule. Within the set of cases to be screened, cases are admitted in increasing order of the predicted failure probability while $k$ times the mean predicted probability among admitted cases stays at or below 5\%. The rule needs no label from the new data. The factor $k \geq 1$ guards against a score that is too optimistic in a new region; it is the 80th percentile, over the other source regions, of the ratio between realised and predicted risk when $k=1$ (\TErFactor{} for R7b, the primary score, and \TErFactorA{} for R7a). The rule was developed by leaving one region out, and Australia, opened by the first test, served as a development region.

The second test used USGS stations whose water temperature had not been read by this or the earlier study when it was designed: \TNewCandidates{} candidates retained by an earlier metadata intake \cite{usgsogc2026}. From the dates of approved daily means alone, without any value, we derived eligibility (at least four complete years), networks (stations in the same eight-digit hydrologic unit, which are the only donors), independence groups (hydrologic units joined by known network connections) and a missingness stratum (stations with incomplete years in at least half of their active years). A second preregistration then fixed four endpoints (N1--N4, Table~\ref{tab:endpoints}): the realised risk under the label-free rule in both populations (N1), the coverage of the calibrated intervals where more than three quarters of the climatological peak is hidden (N2), the bias reduction in the stratum of P1 (N3), and the gain in retention over the best published rule at equal realised risk (N4, descriptive). The definitions, D2*, the interval calibration, the rule, $k$ and all code were frozen; only then were the values requested, and the evaluation ran once. Deployment cases of this test use each station's own gap calendars. Two versions of the corresponding results section were written before the values were read; Section~\ref{sec:transfer} uses the one that matches the outcome reported here.

After both tests had been evaluated, the generation of random numbers was changed so that every seed is a plain integer: a fixed number, or one formed from the position of a station or region in a sorted list, the year and the step of the design. This redrew every mask, calendar draw, shift, simulation and resample in the study. All results, including both tests, were then recomputed with otherwise unchanged definitions, models, rules and code, and the numbers reported here are those of the recomputation. Because the change was made after both outcomes were known, neither test is blind any more, and we report them as evaluations on two additional data sets rather than as confirmatory tests. Text~S6 lists the values of the original evaluations.

\subsection{Downstream experiment}
Stations with at least 20 complete years (\TTrendStations{} stations in \TTrendGroups{} independence groups; 15 years as a sensitivity analysis) were degraded 200 times. Each year became incomplete with the agency's real share of incomplete years, under four scenarios: at random; with the observed decline of missingness over time; with this decline plus the estimated air-temperature effects; and with doubled effects. An incomplete year received one of its real-calendar cases, weighted by the onset importance weight in the state-dependent scenarios. The pipelines were: drop every incomplete year (A0); keep years with at least 90\% warm-season coverage and compute H30 from recorded days (A1); keep years whose longest gap is at most 49 days \cite{johnson2021} (A2); keep all years with the plug-in value (B); keep all years with the D2* median (C1); and keep all years with 20 multiple imputations from D2*, pooled by Rubin's rules with Barnard--Rubin degrees of freedom \cite{rubin1987,barnard1999} (C). Estimands were the ordinary least-squares and Theil--Sen trends of H30, the mean H30, the frequency of years above the complete-record 75th percentile, H30 in regional heat-wave summers (top 10\% of the region's summer air anomaly) and the classification of hot years (above the complete-record 90th percentile). The gain of C over A1 is split into the step from A1 to B, which keeps the years, and the step from B to C, which changes the estimator.

\subsection{Annual 7-day minimum flow}
To see whether the mechanism is specific to temperature, we repeated the design for the annual 7-day minimum flow (7Q) on the USGS daily discharge, using approved values, climatic years from April to March and perennial streams only. Gaps were filled by a ridge regression of log-flow anomalies on up to three neighbours within 200\,km (correlation at least 0.6, complete in the target year), or by interpolating the log anomaly when no neighbour was available. The target-aware estimator was a distribution regression of the log error of the plug-in with the decision-time features used for H30 (as P5a), cross-fitted by leaving out one two-digit hydrologic unit. The tolerance is half the median interannual standard deviation of log 7Q, a relative error of \TFlowTol\%. The USGS practice in SWToolbox skips any climatic year with a missing day \cite{kiang2018}.

\subsection{Preregistration and inference}
The first preregistration was written after the missingness audit and before any estimator comparison. It fixed four primary endpoints with success criteria, the populations, the tolerances and the test region, and controlled the family-wise error over the four endpoints with Holm's method. P1: the reduction of the mean signed H30 error of D2* relative to P0 for air-layer cases whose true peak window is at least half hidden, with 90\% interval coverage between 85\% and 95\%. P2: the retention gain at 5\% realised risk over the best published rule under real calendars, at least 10 points. P3: at least 80\% of held-out regions with realised risk at most 7\% at nominal 5\% and a median of at most 5\%, in both populations, and at most 7\% in Australia. P4: a reduction of at least 20\% in station trend RMSE of C relative to A1, with no worse bias and coverage between 90\% and 97\%. Coverage in P4 refers to the missing-data component $\pm1.96\sqrt{(1+1/M)B}$ around the complete-record slope; this clarification was recorded before the experiment ran. The second preregistration fixed N1--N4; N1--N3 are point-estimate criteria and are judged separately. Both documents and their timelines are in Text~S6, with the list of deviations. Everything else is exploratory. Intervals resample independence groups.

\section{Results}

\subsection{Missing data are only weakly related to heat}
Hot summers did not cause more incomplete years. Within stations, a one-standard-deviation warmer summer raised the odds of a warm-season gap by a factor of \TOrAir{} [\TOrAirLow, \TOrAirHigh] (Figure~\ref{fig:missingness}a). At the daily level, gaps were slightly more likely to begin after warm weeks (odds ratio \TOrOnset{} [\TOrOnsetLow, \TOrOnsetHigh]) and on low-flow days (\TOrLowFlow{} [\TOrLowFlowLow, \TOrLowFlowHigh]), but not after warm water days, and gap length did not depend on the state at onset. The share of incomplete warm seasons in the regionally hottest tenth of summers differed from that in other summers by \THotContrast{} percentage points [\THotContrastLow, \THotContrastHigh] (Figure~\ref{fig:missingness}c). Through the station relation between water and air H30, a complete-years-only climatology would be biased by only \TImpliedBias\,$^{\circ}$C [\TImpliedBiasLow, \TImpliedBiasHigh]. The strongest signal was temporal: the odds of an incomplete warm season fell by a factor of \TOrDecade{} per decade (Figure~\ref{fig:missingness}b). The preregistered gate for state-dependent missingness (odds ratio of at least 1.2) was therefore not met, and state dependence is treated as a stress test below.

\subsection{A recovery law predicts where the plug-in fails, but underpredicts the hardest cases}
The parametric law explained \TLawCellFailEV\% of the variance of failure rates across the 27 cells of gap length, peak overlap and information layer, without fitting any error (Figure~\ref{fig:law}a). At the level of single cases, its Brier skill relative to the base rate was \TLawBrierSkillMech{} and its area under the ROC curve \TLawAucMech{} for controlled gaps (real calendars: \TLawBrierSkillDeploy{} and \TLawAucDeploy), and its expected calibration error was \TLawEceMech{} and \TLawEceDeploy. It underpredicted failure where it matters most. In the five hardest cells, all with most of the peak hidden, it predicted \THardVTwo\% failures against \THardObserved\% observed. The errors there have heavier tails than the law's mixture of normal distributions: \THardBeyond\% fell more than two predicted standard deviations from the predicted mean, against 4.6\% for a normal distribution, the excess kurtosis was \THardKurtosis{} and the variance was \THardVariance\% larger than predicted. Larger residuals in hot periods do not explain the shortfall, because the residual standard deviation inside the training peak windows was \THardHetero{} times its warm-season value. Taking the residual moments from the peak windows (v2p) matched the failure rate of the hardest cells (\THardVTwoP\%) but described the mean error of the cells worse than a constant (explained variance \TVTwoPCellBiasEV).

The law explained less of the mean signed error (\TLawCellBiasEV\%; \TLawCellBiasEVUntruncated\% without truncation) because it overstates the cold bias of neighbour regressions, whose residuals share the heat anomaly with the neighbours. The empirical replay explained \TReplayCellBiasEV\% of the cell variance of the mean error (Figure~\ref{fig:law}b) and predicted an overall bias of \TReplayBiasMech\,$^{\circ}$C against \TObservedBiasMech\,$^{\circ}$C observed, which shows how much of the error the training years contain, but it does not explain it. The law accounts for the difference between information layers: the median $R^2$ of air-temperature reconstructions was \TLawRTwoAir{}, against \TLawRTwoDonor{} for neighbour reconstructions, so the same hidden peak is regressed towards the seasonal cycle much more strongly.

\subsection{Plug-in recovery understates the extreme; target-aware estimation removes most of the bias}
Where the true peak window was at least half hidden and air temperature filled the gap (\TPOneN{} cases, \TPOneGroups{} groups), P0 underestimated H30 by \TPOneBiasPZeroAbs\,$^{\circ}$C on average with a mean absolute error of \TPOneMaePZero\,$^{\circ}$C (Figure~\ref{fig:estimators}a). D2* (P5b, the lowest quantile CRPS: \TCrpsPFiveB, against \TCrpsPFiveA{} for P5a, \TCrpsPThree{} for P3 and \TCrpsPThreeRaw{} for raw residual blocks) reduced the bias to \TPOneBiasCal\,$^{\circ}$C, a reduction of \TPOneReductionCal\% [\TPOneReductionCalLow, \TPOneReductionCalHigh], and the mean absolute error to \TPOneMae\,$^{\circ}$C. Endpoint P1 was not met as registered: the reduction of \TPOneReductionRone\% passed its criterion of 50\%, but the 90\% intervals covered \TPOneCoverageRone\%, just below the registered floor of 85\%. Under the corrected season definitions the uncalibrated intervals covered \TPOneCoverage\%, and less where most of the climatological peak was hidden (Figure~\ref{fig:estimators}c). After calibration, coverage was \TCqrHardMin--\TCqrHardMax\% in every group with a hidden peak and \TCqrEasyMin--\TCqrEasyMax\% in the groups where the peak was visible, where many errors are exactly zero (Table~S19); in the P1 stratum it was \TPOneCoverageCal\% [\TPOneCoverageCalLow, \TPOneCoverageCalHigh]. Because the visible-peak groups over-cover, the intervals are approximate rather than exact. Functional debiasing alone reached \TPOneBiasPTwo\,$^{\circ}$C but gives no interval; air-state simulation reached \TPOneBiasPThree\,$^{\circ}$C and raw residual blocks \TPOneBiasPThreeRaw\,$^{\circ}$C; Air2stream-3 reached \TAirTwoStreamBias\,$^{\circ}$C. Stratified by the climatological peak overlap that a user can see, D2* was nearly unbiased in every stratum (Figure~\ref{fig:estimators}a,b). Removing ridge shrinkage did not change the cold bias (\TShrinkRidge{} and \TShrinkOls\,$^{\circ}$C), which confirms that it arises from regression to the conditional mean.

Some of the error cannot be recovered from visible information. In unusually warm years (true H30 at least 1\,$^{\circ}$C above the training mean) with a mostly hidden peak, D2* still underestimated H30 by \TWarmAnomBiasDTwoAbs\,$^{\circ}$C (P0: \TWarmAnomBiasPZeroAbs\,$^{\circ}$C). Computing H30 from recorded days only was worse than any reconstruction in this situation, underestimating by \TObsDaysBiasHiddenAbs\,$^{\circ}$C when the true peak was at least half hidden.

\subsection{Completeness rules discard mostly usable years}
Under real gap calendars most incomplete years carry little risk: without any screening, \TDeployFailPZero\% of plug-in values and \TDeployFailDTwo\% of D2* values failed, against \TDeployFailObs\% for values computed from recorded days. Apart from the complete-year rule, which uses none of them by definition, the published rules used between \TRetainMin\% and \TRetainMax\% of incomplete years (Table~\ref{tab:rules}). The rules that count days were the most wasteful: the USGS indicator rule used \TRetainUSGS\%, the WMO rule \TRetainWMO\% and the ETCCDI rule \TRetainETCCDI\%. The rules do separate risk, since the years they rejected would have failed more often with the plug-in (\TRejectedFailPZeroUSGS\% for the USGS rule) than the years they kept, but among the rejected years \TRejectedUsableMin--\TRejectedUsableMax\% could have been recovered within tolerance by D2*, for example \TRejectedUsableUSGS\% for the USGS rule and \TRejectedUsableBeaufort\% for the rule that requires two complete peak months. Because D2* without screening already met a 5\% risk, the preregistered retention gain at 5\% was \TPTwoGain{} points and endpoint P2 was not met. At a 2.5\% risk, a warm-season missing-fraction rule with an optimally chosen threshold used \TDeployRetWarmTwoFive\% of years, the target-aware judgment \TDeployRetTargetTwoFive\% and an oracle \TDeployRetOracleTwoFive\% (Figure~\ref{fig:rules}a). With the threshold chosen on the other regions, the warm-season rule used \TLoroFreeWarmRet\% of years, but in the worst held-out region \TLoroFreeWarmWorst\% of them failed. The rule type is nearly efficient; the published thresholds are not, and a free threshold does not transfer reliably between regions.

Screening matters for long warm-season gaps. In the mechanism population, the target-aware judgment with the D2* value used \TMechRetTarget\% of cases at 5\% realised risk (R7a: \TMechRetTargetA\%), against \TMechRetWarm\% for the warm-season completeness rule, \TMechRetTargetDisagree\% for H30 reconstruction disagreement, \TMechRetLaw\% for the recovery law and \TMechRetOracle\% for an oracle (Figure~\ref{fig:rules}b).

Usability depended on the target (Figure~\ref{fig:rules}c,d). Agreement on usability between H30 and other targets was moderate for H7 ($\kappa=\TKappaHSeven$), low for the annual mean (\TKappaMean) and near zero for C30 (\TKappaCThirty; \TKappaCThirtyMech{} under long gaps). Peak-season rules protected H30 but not cold extremes: the rule that requires 80\% of the peak months falsely accepted \TOregonFalseAcceptH\% of years for H30 and \TOregonFalseAcceptC\% for C30. Under long gaps, the count of days above 20\,$^{\circ}$C retained \TDaysOwnScore\% of cases with its own score but only \TDaysHScore\% with the H30 score.

\subsection{Keeping incomplete years improves trends, mostly because the years are kept}
Dropping every incomplete year left an average of \TYearsAZero{} of \TYearsAll{} years and produced station trend errors of \TTrendRmseAZero\,$^{\circ}$C per decade (Figure~\ref{fig:downstream}a). The 90\% warm-season rule had an error of \TTrendRmseAOne, the plug-in with all years \TTrendRmseB{} and multiple imputation from D2* \TTrendRmseC\,$^{\circ}$C per decade: a reduction of \TTrendReduction\% [\TTrendReductionLow, \TTrendReductionHigh] relative to the rule (\TTrendReductionRone\% as registered). The missing-data intervals of C covered the complete-record trend in \TTrendCoverage\% of replicates. Endpoint P4 was met. The results were the same under the four missingness scenarios and with \TTrendStationsFifteen{} stations of at least 15 years (reduction \TTrendReductionFifteen\%).

Most of this gain comes from keeping the years. Moving from the rule to the plug-in with all years achieved \TShareKeep\% [\TShareKeepLow, \TShareKeepHigh] of the reduction from A1 to C, and the estimator the rest. At the \TNStationsRatio{} stations where real calendars reached the peak at least once, the station RMSE ratio of B to A1 had a median of \TRatioBAOne{} (quartiles \TRatioBAOneLow{} and \TRatioBAOneHigh) and that of C to B a median of \TRatioCB{} (\TRatioCBLow{} and \TRatioCBHigh); C was better than B at \TRatioCBBelow\% of these stations (Figure~\ref{fig:downstream}b). At the other \TNStationsBExact{} stations, every transplanted calendar left the peak visible and all pipelines that keep the years were exact. The error grew with the number of incomplete years for every pipeline, and the ordering was the same at every number (Figure~\ref{fig:downstream}c).

On the scale of the trends themselves, these errors are small. The median complete-record trend was \TMedianTrend\,$^{\circ}$C per decade; the RMSE of the 90\% rule was \TRmseAOneShare\% of it and that of C \TRmseCShare\%. Where the complete record detected warming, A1 detected it in \TDetectAOne\% of the replicates and C in \TDetectC\%, with nearly equal interval widths (\TCiWidthAOne{} and \TCiWidthC\,$^{\circ}$C per decade). Dropping every incomplete year detected it in only \TDetectAZero\% and produced false detections in \TFalseDetectAZero\% of the other replicates. The choice between a lenient rule and estimation therefore matters little for detecting warming at single stations with 20--30 years; it matters for the size of the trend, for regional averages and for single years. Among truly hot years that had become incomplete, dropping lost all of them, the 90\% rule missed \THotMissedAOne\%, the plug-in \THotMissedB\% and multiple imputation \THotMissedC\%. The plug-in also biased the mean H30 by \TMeanBiasB\,$^{\circ}$C.

Applied to the \TRealYears{} real incomplete station-years of the archive (\TRealStations{} stations), the label-free target-aware rule admitted \TRealAdmitted\%, the most permissive published rule \TRealBestRule\% and the USGS indicator rule \TRealUSGSRule\%. In Baden-W\"urttemberg a complete-year rule is not feasible: only \TBwCompleteShare\% of the station-years of the \TBwStations{} long-record stations (\TBwGroups{} independence groups) are complete, none of the \TBwIncompleteTwoThousandThree{} stations with data in 2003 recorded that year completely, and in \TBwYearsFewComplete{} of \TBwYears{} years fewer than five stations had a complete year. On the same stations, the regional H30 trend from complete years alone was \TBwAZero\,$^{\circ}$C per decade [\TBwAZeroLow, \TBwAZeroHigh], because it describes mostly the recent part of the record; with incomplete years it was \TBwAOne{} under the 90\% rule, \TBwB{} with the plug-in and \TBwC{} [\TBwCLow, \TBwCHigh] with multiple imputation, and \TBwAOneNoTwoThousandThree, \TBwBNoTwoThousandThree{} and \TBwCNoTwoThousandThree{} without 2003. In a regional anomaly series from station and year effects, 2003 was the warmest year of the record when incomplete years were included and was absent with complete years only (Figure~\ref{fig:downstream}d). No truth is available for these years; the example shows that the treatment of incomplete years decides which period and which summers a regional record can describe.

\subsection{Fixed admission thresholds did not transfer}
Across \TLoroRegions{} held-out source regions, layer-specific thresholds kept realised risk at or below 7\% in \TLoroShareDeploy\% of regions under real calendars (median \TLoroMedianDeploy\%) and in \TLoroShareMech\% under long gaps (median \TLoroMedianMech\%). Because the median exceeded 5\% under long gaps, the source part of endpoint P3 was not met; under the corrected season definitions the share was \TLoroShareMechRtwo\% (median \TLoroMedianMechRtwo\%).

The first test, in Australia, failed. Australian records were harder to recover: \TAuAirShare\% of cases relied on air temperature, the median $R^2$ of the reconstruction was \TAuRTwo{} (source: \TSourceRTwo), and real gap calendars hid a median \TAuMissing\% of the year (source: \TSourceMissing\%). Without screening, \TAuUnscreenedDTwo\% of D2* values failed. The frozen thresholds admitted cases with \TAuRiskDeploy\% realised risk under real calendars [upper bound \TAuRiskDeployHigh\%] and \TAuRiskMech\% under long gaps. D2* reduced the cold bias in the P1 stratum by \TAuPOneReduction\%, but its 90\% intervals covered only \TAuPOneCoverage\%. At equal realised risk, the target-aware judgment used \TAuGain{} percentage points more Australian station-years than the best published rule. From decision-time information alone, the recovery law predicted a plug-in failure rate of \TAuLawPredicted\% in Australia (observed \TAuObservedFail\%), against 3.3\% predicted and 3.6\% observed in the source.

The design of this test had a flaw that we found afterwards: the frozen pipeline used calendar years and a May--September season, so \TAuEdgeBefore\% of the Australian station-years had their H30 window at the calendar-year boundary. With July--June hot years the share fell to \TAuEdgeAfter\%, but Australia remained hard: in these post hoc analyses D2* still failed in \TAuDevUnscreenedDeploy\% of real-calendar cases. A post hoc diagnosis showed why fixed thresholds fail in such a region (Table~S17). The scores ranked cases about as well as in the source, but thresholds set where the base failure rate was below 5\% admitted almost everything, and the same thresholds admit riskier cases where the base rate is higher. This motivated the label-free rule. Developed by leaving one region out, it kept the realised risk at or below 7\% under long gaps in \TErShareKOneMech\% of the \TErRegions{} held-out regions with $k=1$ (worst region \TErWorstKOneMech\%) and in \TErShareMech\% with the factor $k$ (worst \TErWorstMech\%), while using a median \TErMedianRetMech\% of the cases; under real calendars it admitted almost every case and kept the risk at or below 7\% in every region (Table~S12). In Australia, now a development region, it held the risk at \TAuDevRiskDeploy\% under real calendars [upper bound \TAuDevRiskDeployHigh\%] while using \TAuDevRetDeploy\% of the station-years, and at \TAuDevRiskMech\% under long gaps [\TAuDevRiskMechHigh\%] while using \TAuDevRetMech\%.

\input{transfer_versions/transfer_pass.tex}
\label{sec:transfer}

The guarantee holds for the screened set as a whole, not for every part of it: among independence groups with at least 20 admitted cases, the worst reached \TNewWorstGroupDeploy\% under real calendars and \TNewWorstGroupMech\% under long gaps. The D2* failure probability (R7a) gave a similar result (\TNewRiskADeploy\% and \TNewRiskAMech\%), but without the factor $k$ its risk under long gaps would have been \TNewRiskAKOneMech\%. Applying the rule within each hydrologic region instead of to all new stations, or leaving out the targets that share a hydrologic unit with earlier records, changed the risk under long gaps to \TNewRiskPerRegion\% and \TNewRiskNoPrior\%. When $k$ was re-estimated from a few labelled local groups, the risk on the remaining groups stayed at or below 7\% in at least 80\% of repeats under real calendars with two groups and under long gaps with \TNewLocalGroups{} groups (Figure~\ref{fig:transfer}c). Under long gaps the recovery law predicted a plug-in failure rate of \TNewLawMech\% against \TNewObservedMech\% observed, the same underprediction as in the hardest source cells, whereas the peak-window version predicted \TNewLawPMech\%.

\subsection{The same mechanism biases annual low flows}
For the annual 7-day minimum flow at \TFlowSites{} perennial USGS streams (\TFlowYears{} complete climatic years), the plug-in overestimated the minimum when the gap hid the low-flow period, as it underestimates the peak of temperature (Figure~\ref{fig:lowflow}a). Under real gap calendars that hid at least half of the climatological low-flow window, the plug-in 7Q was \TFlowBiasHalfPZero\% too high in the geometric mean, \TFlowBiasInterpPZero\% without a complete neighbour, and 7Q computed from recorded days was \TFlowBiasHalfRecorded\% too high. The SWToolbox practice used none of the incomplete climatic years, although \TFlowUsable\% of them could have been estimated within tolerance; at 5\% realised risk the target-aware judgment used \TFlowRetTarget\% of them and an oracle \TFlowRetOracle\% (Figure~\ref{fig:lowflow}b). Keeping all years reduced the error of station trends in log 7Q by \TFlowTrendReduction\% relative to skipping incomplete years (\TFlowTrendStations{} stations), and \TFlowShareKeep\% of this gain came from keeping the years (Figure~\ref{fig:lowflow}c). The distribution regression improved single values only modestly (failure rate \TFlowFailDTwo\% against \TFlowFailPZero\% for the plug-in) and did not remove the small bias of neighbour fills. For low flows, as for H30 trends, the main gain lies in keeping incomplete years.

\section{Discussion}

\subsection{Why filling and then maximising is biased}
The bias of the plug-in H30 follows from the selection of the peak. The hottest window is, by construction, a period of positive anomaly, and a conditional-mean reconstruction regresses hidden anomalies towards the seasonal cycle in proportion to the share of variance it cannot explain at the window scale. The recovery law makes this quantitative and identifies the variables that matter: the share of the peak window hidden, the persistence of the reconstruction residual relative to the persistence of the anomaly, and the size of typical peak anomalies. These are available before any value is accepted, which is why a law with no fitted error parameters orders gap geometries correctly. Its failure probabilities are too low in the hardest cells, where the errors have heavy tails that a mixture of normal distributions does not reproduce, so they are best read as lower bounds there. For users, the law gives an interpretable reason why the same gap is harmless at a station with good neighbours and harmful at a station that depends on air temperature.

The same argument applies to any annual extreme of a completed series. The minimum of a completed discharge series is biased upwards when the low-flow period is hidden, and computing it from recorded days is worse. The bias is not a property of ridge regression and is only partly removed by process models. Air2stream-3 reduced it, because its anchoring at the last observation preserves some of the anomaly. Estimators that target the statistic directly removed most of it. The remaining cold bias in unusually warm years marks the limit of recovery from visible information: the part of a heat anomaly that neither air temperature nor neighbours record cannot be reconstructed.

\subsection{What the rules cost}
Completeness rules are protective by design, and in our real calendars the years they keep are almost always accurate. The cost lies in what they reject. Day-counting rules reject years for gaps outside the warm season that do not affect a summer extreme, and the published peak-season thresholds are stricter than needed. Rejected years are disproportionately early in the record, because missingness has declined over the decades, so dropping them shortens the effective record, inflates the uncertainty of trends and can change the period that a regional trend describes. The downstream experiment shows that this is the dominant error for trend estimation: most of the gain of estimation over a lenient rule comes from keeping the years, and a better estimator adds a smaller improvement. For single years, such as hot-year counts or values for a regulatory window, the estimator matters more, because the plug-in is biased exactly in the years of interest. The recommendation is therefore to estimate every incomplete year and carry its uncertainty forward, and to use a screening rule only when a single value must be reported. When a rule is needed, its threshold should be set for the target statistic: rules that suit H30 misjudge cold-season statistics, and a score for one target is a poor screen for another.

\subsection{Transfer}
Two preregistered tests bracket what transfers, although neither is blind after the recomputation (Section~\ref{sec:tests}). Thresholds fixed in probability units failed in Australia, where the base failure rate was four times that of the source, although the scores still ranked cases correctly. A calibrated failure probability carries its own expected risk, so a threshold can be set from the screened set's own score distribution without labels, and a factor estimated across regions protects against moderate miscalibration \cite{chow1970,franc2023}. This rule kept the risk at or below 7\% at new USGS stations under their own calendars and under long controlled gaps, and \TNewLocalGroups{} labelled groups were enough to re-estimate the factor; because the evaluation was recomputed after its first outcome was known, the result needs a new independent test. The guarantee is an average over the screened set: stations with many incomplete years and single hydrologic regions came close to 7\% or exceeded it, and one small group was far above it, so the rule should be applied to sets large enough for an average to be meaningful. The second test used new stations in regions that the source represents; a new agency or climate remains untested, and the Australian development results suggest that the factor $k$ matters most there. From decision-time information alone, the recovery law predicted that Australia would be about four times harder than the source and the new stations only moderately harder (\TNewLawPredicted\% against 3.3\% predicted plug-in failures under real calendars), which suggests that it can serve as a warning before thresholds are transferred.

\subsection{Scope}
All errors are evaluated by masking complete years, with gap calendars taken from real incomplete years. Because missingness was only weakly related to heat in this archive, complete years appear to be representative, and the state-dependent scenarios did not change the conclusions; an archive with strong heat-related failures would need its own audit. No natural gap with an independent reference sensor of sufficient accuracy was available: the remaining paired USGS channels differed by 0.2--1.0\,$^{\circ}$C on overlapping days. The intervals of D2* are approximate: calibrated by decision-time group, they cover close to 90\% where the peak is hidden and over-cover where it is visible. The first test used calendar years and a Northern Hemisphere season, which we corrected afterwards; the second used new stations within regions that the source represents, so it does not test a new climate or agency. Both tests were recomputed after their first outcomes were known (Section~\ref{sec:tests}), so the transfer of the label-free rule still needs a blind test. The station trend errors are dominated by a few stations with long records and large errors, so the size of the gain in P4 depends on the random draws: it was 76\% in the first run and \TTrendReductionRone\% in the recomputation, with the same conclusion. The tolerance of 0.5\,$^{\circ}$C and the 5\% risk level are conventions; the conclusions hold at 0.25 and 1.0\,$^{\circ}$C and at 2.5\% and 10\% risk (Tables~S7 and S16). The study used daily means, so regulatory metrics based on daily maxima were not evaluated. Return levels were not estimated because 20--30 years of annual values do not support stable extreme-value fits; methods that correct block maxima for missing data under a random-missingness assumption \cite{simpson2026} could be combined with the multiple imputations used here.

\section{Conclusions}
Incomplete river-temperature years are more useful than completeness rules assume. Filling gaps and then taking the hottest 30 days systematically understates the annual extreme when its window is hidden; a recovery law built from training data alone predicts where this happens, and an estimator aimed at the statistic removes most of the bias and provides approximate intervals. Published completeness rules reject many years that could be recovered, and dropping them degrades station warming trends far more than estimating them does; most of the gain comes from keeping the years rather than from the estimator. Thresholds fixed in probability units did not transfer to a harder region, whereas a label-free threshold based on the expected risk of the screened set held the admitted risk at new stations, in an evaluation that still needs independent confirmation. The same mechanism affects the minimum of completed discharge series. We recommend that archive users estimate the target statistic for every incomplete year with its uncertainty, propagate that uncertainty into trends by multiple imputation, and, where a screen is required, set its threshold from the expected risk of the screened set rather than in fixed probability units.

\section*{Open Research}
Provider data are available from the agencies under their terms; the acquisition routes and station lists are archived with the analysis code. Code, configuration, both preregistrations and derived results are in the project repository; a persistent archive with a DOI will be created on acceptance and cited here. PENDING: repository DOI.

\acknowledgments
PENDING AUTHOR CONFIRMATION: funding and acknowledgments.

\begin{figure}
\centering
\includegraphics[width=\textwidth]{figures/tr_f1_missingness.pdf}
\caption{Missingness in the archive. (a) Odds ratios with 95\% intervals for incompleteness and gap onset; the dashed line marks the preregistered gate. (b) Share of fully operational station-years with a May--September gap by decade. (c) Share of incomplete warm seasons in each region's hottest 10\% of summers and in other summers.}
\label{fig:missingness}
\end{figure}

\begin{figure}
\centering
\includegraphics[width=\textwidth]{figures/tr_f2_recovery_law.pdf}
\caption{The recovery law. (a) Predicted and observed failure rates of the plug-in H30 for each gap length, peak overlap and information layer (parametric law). (b) Predicted and observed mean signed error from the training-year replay, an empirical comparator. (c) Reliability of the parametric law for controlled gaps and real calendars.}
\label{fig:law}
\end{figure}

\begin{figure}
\centering
\includegraphics[width=\textwidth]{figures/tr_f3_estimators.pdf}
\caption{Bias and coverage of H30 estimators for controlled gaps, by the share of the climatological peak window hidden. (a) Air-temperature layer. (b) Neighbour layer. (c) Coverage of 90\% intervals before and after calibration by decision-time group; the band marks 85--95\%.}
\label{fig:estimators}
\end{figure}

\begin{figure}
\centering
\includegraphics[width=\textwidth]{figures/tr_f4_rules.pdf}
\caption{Applicability. (a) Risk--retention curves under real gap calendars; numbered points are published rules with values computed from recorded days. (b) The same under controlled 30--90 day gaps. (c) Retention at 5\% realised risk for six targets under controlled gaps. (d) Agreement on usability between targets.}
\label{fig:rules}
\end{figure}

\begin{figure}
\centering
\includegraphics[width=\textwidth]{figures/tr_f5_transfer.pdf}
\caption{Transfer of admission thresholds. (a, b) Realised risk among admitted cases for each held-out source region under the label-free rule with $k=1$ and with the factor $k$; below, the fixed thresholds in the first test (Australia), the label-free rule in Australia after it was opened (post hoc), and the label-free rule in the second test on new USGS stations, both tests as recomputed (Section~\ref{sec:tests}), with 95\% intervals over independence groups. (c) Share of repeats with realised risk at most 7\% on the remaining new stations when $k$ is re-estimated from a few labelled local groups.}
\label{fig:transfer}
\end{figure}

\begin{figure}
\centering
\includegraphics[width=\textwidth]{figures/tr_f6_downstream.pdf}
\caption{Downstream consequences. (a) RMSE of station H30 trends relative to the complete record under four missingness scenarios. (b) Station RMSE ratios for keeping the years (plug-in with all years against the 90\% rule) and for the estimator (multiple imputation against the plug-in), fitted state-dependent scenario; bars show medians and quartiles. (c) Trend RMSE by the number of years made incomplete. (d) Regional H30 anomaly of the same 19 Baden-W\"urttemberg stations from station and year effects, with complete years only and with incomplete years estimated by D2*; real archive, no truth.}
\label{fig:downstream}
\end{figure}

\begin{figure}
\centering
\includegraphics[width=\textwidth]{figures/tr_f7_low_flow.pdf}
\caption{Annual 7-day minimum flow (7Q) at perennial USGS streams. (a) Mean relative error of 7Q under real gap calendars by the share of the climatological low-flow window hidden. (b) Risk--retention curves under real calendars; grey points are fixed rules with values from recorded days, including the SWToolbox practice of skipping incomplete years. (c) RMSE of station trends in log 7Q.}
\label{fig:lowflow}
\end{figure}

\input{tr_tables_main.tex}

\bibliography{references_tr}

\end{document}
